\chapter{Conclusion and Future Scope}

\section{Conclusion}
The Vigil-Core project successfully demonstrates the feasibility and immense value of an integrated, full-stack digital forensics platform. By consolidating disparate backend tools—specifically \texttt{cyart\_malware} and \texttt{NetForensicX}—into a cohesive pipeline, the platform significantly reduces the manual overhead traditionally associated with incident response. 

The dedicated efforts in backend testing ensured robust Linux compatibility and the reliable processing of massive data sets, overcoming initial memory limitations through chunked processing. The integration of YARA rules and fuzzy hashing provided a highly accurate, resilient malware detection engine. Concurrently, the strategic architectural design of the frontend resulted in a stable, low-latency UI that presents complex forensic artifacts as clear, actionable intelligence. Ultimately, Vigil-Core delivers on its objective: empowering security analysts with a manager-ready, automated forensics platform.

\section{Future Scope}
While the current iteration of Vigil-Core is highly capable, several avenues remain for future enhancement:
\begin{itemize}
    \item \textbf{Dynamic YARA Rule Management:} A priority for the next iteration is to allow analysts to dynamically upload, modify, and test new YARA rules directly through the UI, without requiring backend code changes. This feature was planned at the architectural level but deferred due to time constraints.
    \item \textbf{Cloud-Based Sandboxing:} Integrating a fully isolated, cloud-based dynamic analysis environment to safely detonate and monitor highly evasive zero-day threats.
    \item \textbf{AI-Powered Threat Intelligence:} Implementing machine learning models to analyze historical scan data, providing predictive insights and automated threat categorization based on behavioral patterns.
\end{itemize}
